% The final section: complements to the main theorem (label prefix cmp; subsection labels mai). Continued in part 2.
\section{Complements}
\label{sec:cmp}

This section records the consequences of $\Dn=1$ (\cref{thm:gam-D}) for the natural form and proves universality in flat form; none of it is used in the proofs of Theorem~A and Corollary~B. The conventions of \cref{sec:mai} are kept.

\subsection{Consequences of \texorpdfstring{$\Dn=1$}{D=1} for the natural form}
\label{sec:mai-nat}

Let $\hat\Qr(u,\zh)=\Lambda_s\Qr(u,s^4\zh)$ be the hatted rows of \cref{def:set-hat}. Its constant term is $\hat\Qr(u,0)=\Lambda_s\Psi_{\mathrm{cf}}=\Psi(u)$ (\cref{lem:mai-Psi}~(ii)), and $\tilde R(\zh;u)\coloneqq\Psi(u)^{-1}\hat\Qr(u,\zh)=\Psi(u)^T\eta^{-1}\hat\Qr(u,\zh)=\Psi_{\mathrm{cf}}^{-1}\Cexp(\Qr)(s^4\zh)$ is the conifold expansion of the matrix $\tilde R$ of \cref{def:mai-Fcf}, in the variable $\zh$.

\begin{lem}[the hatted gauge]\genusstatus{AI-proposed, unchecked}
\label{lem:mai-Ghat}
Write the graded canonical gauge as $\Pc=\sum_kz^k\sum_bp_{k,b}(u)E_{b+k,b}$ with $p_{k,b}\in\bbQ\fps{u}$ (\cref{thm:gau-U}, \cref{prop:gau-graded}, \cref{def:gau-can}), and put $\Ghat(u,\zh)\coloneqq\Lambda_s\Pc(u,s^4\zh)\Lambda_\varsigma^{-1}$. Then
\[
\Ghat_{ab}=\tau^{(1-\nu_b)/4}\sum_k\zh^k\,u^{\mu(b,k)}\,p_{k,b}(u)\,\delta_{a,b+k},\qquad\mu(b,k)\coloneqq\frac{4k+\nu_b-\nu_{b+k}}{5},
\]
where $\mu(b,k)$ is an integer with $\mu(b,k)\geq\max(0,(4k-4)/5)$, and $\mu(b,k)=0$ only for $k=0$ and for $(b,k)=(1,1)$, the entry $(2,1)$. Hence $\Ghat$ is a power series in $u$ whose coefficient of $u^n$ is a polynomial in $\zh$ of degree at most $\lfloor(5n+4)/4\rfloor$.
\end{lem}

\begin{proof}
Since $\varsigma=s\tau^{-1/4}$ (\cref{def:set-tau}), the entry $(b+k,b)$ of $\Lambda_s\Pc(u,s^4\zh)\Lambda_\varsigma^{-1}$ is $p_{k,b}\zh^k\tau^{(1-\nu_b)/4}$ times $s$ to the power $(1-\nu_{b+k})+4k-(1-\nu_b)=4k+\nu_b-\nu_{b+k}$. As $\nu_a\equiv-a\bmod5$, this exponent is $\equiv5k\equiv0\bmod5$, so it is $5\mu(b,k)$ and $s^{5\mu(b,k)}=u^{\mu(b,k)}$. The $\nu_a$ take values in $\{-1,\dots,3\}$, so $\nu_b-\nu_{b+k}\geq-4$ and $\mu(b,k)\geq(4k-4)/5$; for $k=0$ the exponent is $0$, and for $k\geq1$ it is $\geq0$. For $k\geq2$ it is positive; for $k=1$ it vanishes iff $\nu_{b+1}-\nu_b=4$, which happens only for $b=1$. At a given $u^n$ only the $k$ with $(4k-4)/5\leq n$ occur.
\end{proof}

\begin{cor}[the natural form after $\Dn=1$]\genusstatus{AI-proposed, unchecked}
\label{cor:mai-natfact}
In every conifold frame of \cref{def:set-frames}:
\begin{enumerate}
\item[(a)] $\hat\Qr(u,\zh)=\Ghat(u,\zh)\,\Rhat^{\mathrm m}(\tau(u)\zh)$. In particular $\hat\Qr$ is a power series in $u$ whose coefficients are power series in $\zh$, and $\Ghat$ is a power series in $u$ whose coefficients are polynomials in $\zh$ (\cref{lem:mai-Ghat}).
\item[(b)] $\hat\Qr(0,\zh)=\Rhat^{\mathrm m}(\zh)$: the conifold symbol of the formal quintic is that of the model. Hence $\tilde R(\zh;0)=\Mo(\zh)$ and $\Psi(0)=\Psi_0$.
\item[(c)] $\tilde R(\zh;u)=G(u,\zh)\,\Mo(\tau(u)\zh)$, with $\Mo$ on the right, where
\[
G\coloneqq\Psi(u)^{-1}\,\Ghat\,\Psi_0=\Psi_0^{-1}\,\Ghat(u,0)^{-1}\Ghat(u,\zh)\,\Psi_0 .
\]
Moreover $G(u,0)=G(0,\zh)=1$, the coefficient of every $u^n$ in $G$ is a polynomial in $\zh$, and $G(-\zh)^TG(\zh)=1$. The factor $\Mo(\tau\zh)$ is the same in every conifold frame; $G$ is built from the canonical gauge of the frame and may depend on it.
\end{enumerate}
\end{cor}

\begin{proof}
(a) By \cref{thm:gau-R0}~(iv), \cref{cor:gau-colours} and $\Dn=1$, $\Qr=\Pc\Qr^{\mathrm m}$, and $\Qr^{\mathrm m}=\Lambda_\varsigma^{-1}\Rhat^{\mathrm m}(z/\varsigma^4)$ with $z/\varsigma^4=\tau\zh$ (\cref{prop:set-modrows}, \cref{def:set-tau}, \cref{lem:set-canon}; \cref{lem:flh-inputs}~(i)). Multiplying by $\Lambda_s$ and inserting $\Lambda_\varsigma^{-1}\Lambda_\varsigma$ gives $\hat\Qr=\Ghat\,\Rhat^{\mathrm m}(\tau\zh)$. The coefficients of $\Rhat^{\mathrm m}(\tau\zh)$ lie in $\bbQ(\zeta)\fps{u}$, so the polynomiality of $\Ghat$ (\cref{lem:mai-Ghat}) makes $\hat\Qr$ a power series in $u$ with coefficients in $\bbQ(\zeta)\fps{\zh}$; in particular $\hat\Qr$ can be evaluated at $u=0$.

(b) At $\zh=0$, (a) gives $\Psi(u)=\Ghat(u,0)\Psi_0$, since $\Rhat^{\mathrm m}(0)=\Psi_0$ ($\Fs_b(0)=\varepsilon_b$, \cref{prop:set-modrows}); with \cref{lem:mai-Psi}, $\Ghat(0,0)=\Psi(0)\Psi_0^{-1}=1$. By \cref{lem:mai-Ghat} the only entry of $\Ghat$ at $u^0$ off the diagonal is the entry $(2,1)$, with coefficient $\zh\,p_{1,1}(0)$, and $p_{1,1}(0)=0$ by \cref{lem:gau-R1}: column $1$ of $\Pc(0,z)$ is $(\Pc)_{11}(0,z)e_1$. So $\Ghat(0,\zh)=1$, and $\tau(0)=1$ gives $\hat\Qr(0,\zh)=\Rhat^{\mathrm m}(\zh)$. Then $\tilde R(\zh;0)=\Psi_0^{-1}\Rhat^{\mathrm m}(\zh)=\Mo(\zh)$ by \cref{def:set-M0}.

(c) $\tilde R=\Psi(u)^{-1}\Ghat\,\Rhat^{\mathrm m}(\tau\zh)=[\Psi(u)^{-1}\Ghat\,\Psi_0][\Psi_0^{-1}\Rhat^{\mathrm m}(\tau\zh)]=G\,\Mo(\tau\zh)$. The second expression for $G$ follows from $\Psi(u)=\Ghat(u,0)\Psi_0$; $\Ghat(u,0)=\diag(\tau^{(1-\nu_b)/4}p_{0,b})$ is invertible over $\bbQ\fps{u}$, the $p_{0,b}$ being units (\cref{thm:gau-U}), so every $u^n$-coefficient of $G$ is a polynomial in $\zh$. $G(u,0)=1$ is clear, and $G(0,\zh)=1$ is (b). For the symplecticity, \cref{lem:flh-inputs}~(iii) with $\Dn=1$ and \cref{lem:set-M0sympl} give $\Qr(-z)^T\eta^{-1}\Qr(z)=1$, hence $\hat\Qr(-\zh)^T\eta^{-1}\hat\Qr(\zh)=1$ because $\Lambda_s\eta^{-1}\Lambda_s=\eta^{-1}$, and $\tilde R(-\zh)^T\tilde R(\zh)=1$ because $\Psi(u)^{-T}\Psi(u)^{-1}=\eta^{-1}$ (\cref{lem:mai-Psi}). Therefore $G(-\zh)^TG(\zh)=\Mo(-\tau\zh)^{-T}\Mo(\tau\zh)^{-1}=1$ by \cref{lem:set-M0sympl}. Neither $\Mo$ nor $\tau$ involves the frame.
\end{proof}

Only the symplecticity in (c) sees $\Dn$, through the even part of $\log\Dn_0$: for a scalar $\Dn=\cn(z)\in1+z^5\bbQ\fps{z^5}$ the proof gives the rest with $\Pc\cn(z)^{-1}$ in place of $\Pc$, as $z^5=u^4\zh^5$ (\cref{rmk:mai-stored}).

\begin{cor}[the actual gauge is canonical]\genusstatus{AI-proposed, unchecked}
\label{cor:mai-actual}
Let $P^{\mathrm{act}}\coloneqq\Qr\,(\Qr^{\mathrm m})^{-1}$, the matrix series that carries the model columns at $v(u)$ to the rows of the formal quintic, both with their exponentials stripped (\cref{thm:gau-R0}). Then $P^{\mathrm{act}}=\Pc$. In particular $P^{\mathrm{act}}$ is a graded regular gauge with coefficients in $\bbQ\fps{u}$, $P^{\mathrm{act}}_{11}(0,z)=1$, and $P^{\mathrm{act}}\YC(v(u))=\YCQ$.
\end{cor}

\begin{proof}
$\Qr^{\mathrm m}$ is invertible because $\Qr^{\mathrm m}(0)=\Psi^{\mathrm m}$ is (\cref{lem:flh-inputs}~(i)). By \cref{thm:gau-R0}~(iv) with \cref{cor:gau-colours}, $\Qr\Dn=\Pc\Qr^{\mathrm m}$, and $\Dn=1$ by \cref{thm:gam-D}; so $P^{\mathrm{act}}=\Pc$. The remaining assertions are \cref{def:gau-can} and \cref{thm:gau-G}.
\end{proof}

With \cref{thm:gam-D} this is Theorem~C (\cref{thm:int-D}); it is not circular, as the rows are normalised at $u=1$ and the canonical gauge at $u=0$.

\subsection{Universality}
\label{sec:mai-univ}

The flat half of the proof of Theorem~A, \cref{cor:red-can} with \cref{thm:mod-MLK}, uses the formal quintic only through its cyclic system and the canonical gauge, and applies to every self-dual cyclic system with one simple pole whose product is that of a pulled-back model; the comparison (\cref{thm:flh-main}) and the normalisation (\cref{thm:gam-D}) are specific to the formal quintic. The systems are over $\bbQ$ (\cref{def:gau-onepole}), and for a non-zero series $y\in u\bbQ\fps{u}$ (a local letter), $\tilde A_y$ denotes the model pulled back by $y$, with coefficients $(y',-y'/y,y',-y',-y')$ in $\bbQ\lser{u}$ and product $-y'^5/y$ (\cref{prop:gau-equalprod}); it has the one-pole shape if and only if $y'(0)\neq0$.

\begin{lem}[the pole order]\genusstatus{AI-proposed, unchecked}
\label{lem:mai-pole}
Let $A=\sum_ab_aE_{a,a+1}$ with $b_a\in\bbQ\lser{u}$ non-zero, $\Bprod\coloneqq\prod_ab_a$, and let $y\in u\bbQ\fps{u}$ be non-zero of order $n\geq1$.
\begin{enumerate}
\item If the gauge equation $z\,\partial_uP=AP-P\tilde A_y$ has a solution whose $z^0$-coefficient $P_0$ is non-zero, then $\Bprod=-y'^5/y$.
\item In that case $\ord_0\Bprod=4n-5\in\{-1,3,7,\dots\}$. If $\Bprod$ has a pole at $u=0$, then $n=1$, $\ord_0\Bprod=-1$ and $y'(0)^4=-\Res_{u=0}\Bprod$.
\end{enumerate}
\end{lem}

\begin{proof}
(i) At $z^0$ the gauge equation reads $AP_0=P_0\tilde A_y$, hence $A^5P_0=P_0\tilde A_y^5$, with $A^5=\Bprod\cdot1$ and $\tilde A_y^5=-(y'^5/y)\cdot1$; so $(\Bprod+y'^5/y)P_0=0$ over the field $\bbQ\lser{u}$, and $P_0\neq0$ gives the claim (the necessity clause of \cref{prop:gau-equalprod}). (ii) $y'^5/y$ has order $5(n-1)-n=4n-5$; a pole forces $n=1$, and then $y=y'(0)u+O(u^2)$ and $y'^5/y=y'(0)^4u^{-1}+O(1)$, so $\Res_{u=0}\Bprod=-y'(0)^4$.
\end{proof}

\begin{cor}[universality, flat form]\genusstatus{AI-proposed, unchecked}
\label{cor:mai-univ}
Let $v_A\in u\bbQ\fps{u}$ with $v_A'(0)\neq0$, and let $A$ be a self-dual one-pole system over $\bbQ\fps u$ (\cref{def:gau-onepole}) with product $-v_A'^5/v_A$ and residue $r$. Then $r=-1/\lambda^2$ with $\lambda\coloneqq-b_0(0)b_3(0)/v_A'(0)^2\in\bbQ^\times$ (a letter local to this corollary and its proof), because $b_1=\Bprod/(b_0^2b_3^2)$ by self-duality and $(u\Bprod)(0)=-v_A'(0)^4$ (\cref{lem:mai-pole}~(ii) with $y=v_A$). Let $P^A_{\mathrm{can}}$ be the canonical gauge from $\tilde A_{v_A}$ to $A$, that is, the regular gauge with $(P^A_{\mathrm{can}})_{11}(0,z)=1$, which exists, is unique, is graded and has rational coefficients, by the argument of \cref{def:gau-can}: a rational graded regular gauge $P$ exists (\cref{prop:gau-graded}~(i)), $P_{11}(0,z)$ is a unit of $\bbQ\fps{z^5}$ (\cref{prop:gau-graded}~(iii)), so $P\,P_{11}(0,z)^{-1}$ is a graded regular gauge (\cref{prop:gau-graded}~(ii)), and it is the only regular gauge with this normalisation and has rational coefficients (\cref{thm:gau-U}~(b), (d)); and let $P_A\coloneqq\lambda^{-1}P^A_{\mathrm{can}}$. Then for every $g\geq2$:
\begin{enumerate}
\item the term of every leg-free stable graph of genus $g$ with an edge in the flat sum of $P_A.\Omega^{\mathrm{mod}}_{v_A(u)}$ lies in $\bbQ\fps{u}$;
\item the term of the vertex graph lies in $c_g(T_A/\lambda)^{2-2g}+\bbQ\fps{u}$, where $T_A\coloneqq\int_0^ub_0\,\dd u$.
\end{enumerate}
\end{cor}

For the formal quintic in a conifold frame, $A$ is the system of \cref{prop:gau-quintic}, $v_A=v(u)$, $r=-1$, $\lambda=1$ ($b_0(0)=1$, $b_3(0)=-1$, $v'(0)=1$), $P_A=\Pc$ and $T_A=\tc$ (\cref{lem:gau-vanq}~(d)), and the corollary is \cref{cor:red-can} with \cref{thm:mod-MLK}. The sign of $\lambda$ is a convention: the proof uses only $\lambda^2=-1/r$, and the flat sums of $P$ and of $-P$ agree term by term, since $R$, $E_P$ and $\tilde q_T$ are unchanged, $d$ changes sign, and the vertex tensors of \cref{def:red-flatsum} change by $(-1)^{n_v}$ with $\sum_vn_v$ even.

\begin{proof}
For the pair $(A,\tilde A_{v_A})$ the residue of $\tilde A_{v_A}$ is $\tilde r=-1$. By \cref{lem:gau-R1}, column $1$ of $P(0,z)$ is $P_{11}(0,z)e_1$ and $rP_{22}(0,z)=\tilde rP_{11}(0,z)$ for every regular gauge $P$. The proof of \cref{prop:gau-symp} uses the residues only in its last step, the evaluation of the entry $(2,1)$ at $u=0$ by \cref{lem:gau-R1}; with the general form of that lemma it gives $P(-z)^T\eta P(z)=(\tilde r/r)\,P_{11}(0,z)P_{11}(0,-z)\,\eta$. For $P_A$ one has $P_{11}(0,z)=\lambda^{-1}$ and $\tilde r/r=\lambda^2$, so $P_A$ is symplectic. By \cref{thm:gau-G}~(c), $P^A_{\mathrm{can}}\YC(v_A(u))=Y^{\mathrm C}_A$, since $\YC(v_A(u))$ is the vanishing solution of $\tilde A_{v_A}$ (\cref{lem:gau-vanmod}, the chain rule, \cref{lem:gau-vanshape}, whose last claim holds for every $v\in u\bbQ\fps u$, and \cref{lem:gau-uniq}, as in the proof of \cref{thm:gau-G}); so $P_A\YC(v_A(u))=\lambda^{-1}Y^{\mathrm C}_A$ has no positive power of $z$ (\cref{lem:gau-vanq}), which is (A), and $T_{P_A}=T_A/\lambda$ by \cref{lem:gau-vanq}~(d). $P_A$ is graded with coefficients in $\bbQ\fps{u}$ and invertible constant term (\cref{thm:gau-U}, \cref{prop:gau-graded}), so it is a dressing (\cref{def:tra-dressing}) with (Sp) and (A). \Cref{thm:red-main} with \cref{thm:mod-MLK} gives (i) and (ii).
\end{proof}

\begin{rmk}[non-gauge dressings]
\label{rmk:mai-annih}
\Cref{thm:red-main} does not use the gauge equation: with \cref{thm:mod-MLK} it gives, for every dressing $P$ with (Sp) and (A) (\cref{def:tra-dressing}) and every $g\geq2$, regular edge graphs in the flat sum of $P.\Omega^{\mathrm{mod}}_{v(u)}$ and the vertex graph in $c_gT_P^{2-2g}+\bbQ\fps{u}$. Only the flat form is claimed: a hatted statement would need an identity like \cref{lem:flh-inputs}~(ii), available only for the formal quintic.
\end{rmk}

\begin{rmk}[a non-canonical normalisation]
\label{rmk:mai-stored}
Let $P^\sharp\coloneqq\Pc\,\cn(z)$, $\cn\in1+z^5\bbQ\fps{z^5}$, be the regular gauge of the other normalisation of \cref{sec:sgn-norm}, and write $\cn(z)=1+m_5z^5+m_{10}z^{10}+\cdots$ ($P^\sharp$ is a local letter). By \cref{thm:gau-U}, $P^\sharp$ is a graded regular gauge. By \cref{prop:gau-symp}, $P^\sharp(-z)^T\eta P^\sharp(z)=\cn(z)\cn(-z)\eta$ with $\cn(z)\cn(-z)=1-(m_5^2-2m_{10})z^{10}+\cdots$: $P^\sharp$ is symplectic modulo $z^{10}$, and symplecticity constrains only the even part of $\log\cn$. By \cref{thm:gau-G}, $P^\sharp\YC=\cn(z)\YCQ$, whose part of positive degree in $z$ begins with $m_5(z^5\YCQ)_{+}=-m_5(z^5e_1+\tc z^4e_0+\cdots)$; by \cref{thm:gau-G}~(b), $P^\sharp$ satisfies (A) only if $\cn=1$. So the gauge equation does not imply (A), by the freedom $\cn(z^5)$ of \cref{thm:gau-U}; (G) fixes $\cn$ altogether, in particular the odd part of $\log\cn$, which symplecticity leaves free. The value of $m_5$ is known only by computation.
\end{rmk}

